\documentclass[10pt,a4paper]{amsart}
\usepackage[margin=19mm]{geometry}
\usepackage{amsmath,amssymb,mathtools}
\usepackage[scr=rsfs,cal=euler]{mathalfa}
\usepackage{xfrac}
\usepackage[unicode,hidelinks,bookmarks=false]{hyperref}
\newcommand{\ie}{i.e.\ }
\newcommand{\cf}{cf.\ }
\newcommand{\resp}{respectively\ }
\newcommand{\Subsection}{\subsection}
\providecommand{\nobreakdash}{\nobreak\hbox{-}}
\setlength{\emergencystretch}{2em}
\multlinegap0pt
\usepackage{bm}
\usepackage{bbm}
\usepackage{dsfont}
\usepackage{xspace}
\usepackage{cancel}
\usepackage{mathtools}
\usepackage{amsthm}
\makeatletter
\@ifundefined{definition}{\newtheorem{definition}{Definition}}{}
\@ifundefined{claim}{\newtheorem{claim}[definition]{Claim}}{}
\makeatother
\makeatletter
\expandafter\ifx\csname pf\endcsname\relax
\newenvironment{pf}{\noindent{\bf Proof}}
\fi
\makeatother
\title[On a refinement of the Ahlswede--Katona Theorem]{On a refinement of the Ahlswede--Katona Theorem}
\author{Jianfeng Hou, Xizhi Liu, Yixiao Zhang}
\date{Source: arXiv:2510.19681v2 (CC BY 4.0)}
\begin{document}
\maketitle

\noindent\emph{Source and licence.} On a refinement of the Ahlswede--Katona Theorem. Version arXiv:2510.19681v2 (CC BY 4.0), distributed under the Creative Commons Attribution 4.0 International licence, as recorded in the arXiv licence metadata for this exact version and shown on the abstract page https://arxiv.org/abs/2510.19681v2. Full rights notice: licenses/arxiv-2510.19681-CC-BY-4.0.txt.\par\smallskip

\noindent\textit{Editorial scope.} Counted unit: the Claim whose source label is \texttt{THM-bi-degree-du1} in the article (source0.tex lines 614--616, source label \texttt{THM-bi-degree-du1}), delivered together with the article's own complete original proof of it (source0.tex lines 618--685, one \texttt{proof} environment of 68 lines). No mathematics is written, repaired or re-derived: statement and proof are the article's own text line for line. Required context retained uncounted: CLAIM:dw1>=du1 (lines 515--517); equ:ZG-and-ZG1 (lines 588--594); equ:induction-ZG1 (lines 605--612). Every definition, display and label of those lines is retained unchanged. Omitted from this package and not redistributed: 1--514, 518--587, 595--604, 613--613, 617--617, 686--1793. Nothing inside a retained excerpt is omitted or altered: every retained body is delivered exactly as the article's lines are written, so no omission receipt of any kind accompanies this package. Citations: the retained excerpts carry no citation command at all, so no bibliography entry of the article is transcribed and no reference is redistributed. Reference adaptation: none was needed and none was made; every reference inside the delivered unit resolves inside it, so no unresolved reference or citation is produced. Printed slips kept literal: no printed word, symbol, hypothesis or number of the counted statement or of its proof was altered. Layout and class: the article's own class is \texttt{article} with packages including epsf, epsfig, amsfonts, amsgen, amsmath, amstext, amsbsy, amsopn, amsthm, cases, listings, color, ebezier, eepic; the wrapper is the lane's pinned \texttt{amsart} editorial wrapper at 10pt on A4, which restates the theorem-like environments the retained text needs and adds the article's own macro definitions that the retained text needs (none), with extra wrapper packages (bm, bbm, dsfont, xspace, cancel, mathtools). The delivered numbering is the unit's own. Assets: no retained span contains a figure environment, a graphics-inclusion command, a PDF-page inclusion, a table or a TikZ picture; no image file is copied into this package, the package declares no asset dependency and no image file is redistributed with it. Licence: version arXiv:2510.19681v2 of this article is distributed under the Creative Commons Attribution 4.0 International licence, as recorded in the arXiv licence metadata for this exact version and shown on the abstract page https://arxiv.org/abs/2510.19681v2. A full rights notice is delivered at licenses/arxiv-2510.19681-CC-BY-4.0.txt. Characters: every retained excerpt is pure ASCII, so the delivered engine, XeLaTeX with its default Latin Modern text font, reports no missing character.
    \begin{claim}\label{CLAIM:dw1>=du1}
        We may assume that $\Delta_{G, \mathrm{right}} \ge \Delta_{G, \mathrm{left}}$.  
    \end{claim}

    \begin{align}\label{equ:ZG-and-ZG1}
        Z_{1}(G) 
        & = \sum_{i=1}^{r}d_{G}(u_i)^2 + \sum_{j=1}^{s} d_{G}(w_j)^2 \notag \\ 
        & =\sum_{i=1}^{d_{G}(w_1)}d_{G}(u_i)^2 + \sum_{j=1}^{d_{G}(u_1)} d_{G}(w_j)^2 \notag \\
        & =\sum_{i=1}^{d_{G}(w_1)} \big( (d_{G}(u_i)-1)^2  + 2d_{G}(u_i) - 1 \big)  + \sum_{j=2}^{d_{G}(u_1)} d_{G}(w_j)^2 +  d_{G}(w_1)^2 \notag \\
        &= Z_{1}(G_1)+  2m  +d(w_1)^2 - d(w_1).
    \end{align}

    \begin{align}\label{equ:induction-ZG1}
        Z_{1}(G_1) \le 
        \begin{cases}
            Z_{1}\big( B_1(\hat{r}, \hat{s}, \hat{m}, \hat{\ell}, \hat{k}) \big), &\quad\text{if}\quad \hat{m} \le \hat{r} \hat{k} ~\text{and}~ \hat{k} + \hat{\ell} \le \hat{r}, \\[0.5em]
            Z_{1}\big( B_2(\hat{r}, \hat{s}, \hat{m}, \hat{\ell}, \hat{k}) \big), &\quad\text{if}\quad \hat{m} \le \hat{r} \hat{k} ~\text{and}~ \hat{k} + \hat{\ell} > \hat{r}, \\[0.5em]
            Z_{1}\big( B(\hat{r}, \hat{s}, \hat{m}) \big), &\quad\text{if}\quad \hat{m} \ge \hat{r} \hat{k}. 
        \end{cases}
    \end{align}

    \begin{claim}\label{THM-bi-degree-du1}
        We may assume that $\Delta_{G, \mathrm{left}} = k$. 
    \end{claim}

    \begin{proof}[Proof of Claim~\ref{THM-bi-degree-du1}]
        Suppose that $\hat{m} \le \hat{r} \hat{k}$. 
        Then let $G^{\ast}$ denote the graph obtained from $G$ by replacing $G_1$ with $B_1(\hat{r}, \hat{s}, \hat{\ell}, \hat{k}, \hat{m})$ or $B_2(\hat{r}, \hat{s}, \hat{\ell}, \hat{k}, \hat{m})$ (depending on whether $\hat{k} + \hat{\ell} \le \hat{r}$).
        It follows from the definition of $B_1(\cdot)$ and $B_{2}(\cdot)$ that $G^{\ast} \in \mathcal{B}(r, s, \ell, k, m)$. 
        Moreover, it follows from~\eqref{equ:ZG-and-ZG1} and~\eqref{equ:induction-ZG1} that $Z_{1}(G^{\ast}) \ge Z_{1}(G)$. 
        Therefore, $G^{\ast} \in \Phi(r,s,\ell,k,m)$. 
        Note that $\Delta_{G^{\ast}, \mathrm{left}} = \hat{k} +1 =k$. 
        We are done. 

        Now suppose that $\hat{r} \hat{k} < \hat{m} = m-d_{G}(w_1)$. 
        Then it follows that 
        \begin{align*}
            d_{G}(w_1)
            = \hat{r}
            < \frac{m}{\hat{k}+1}
            \le \frac{rk}{k}
            = r, 
        \end{align*}
        which implies that $d_{G}(w_1) \le r-1$. 
        
        Let $G^{\ast}$ be the graph obtained from $G$ by replacing $G_1$ with $B(\hat{r}, \hat{s}, \hat{m})$. 
        Then $|G^{\ast}| = |G|$. 
        Since $\hat{m} > \hat{r}\hat{k}$, the definition of $B(\hat{r}, \hat{s}, \hat{m})$ implies that there are at least $\hat{k}$ vertices in $W_1$ with degree $\hat{r} = d_{G}(w_1) \ge \ell$. 
        Together with the vertex $w_1$, there are at least $\hat{k} + 1 = k$ vertices in $W$ adjacent to every vertex in $I$. 
        Consequently, $\delta_{G^{\ast}}(I) \ge k$, which implies $G^{\ast} \in \mathcal{B}(r,s,\ell,k,m)$. 
        Moreover, it follows from~\eqref{equ:ZG-and-ZG1} and~\eqref{equ:induction-ZG1} that $Z_{1}(G^{\ast}) \ge Z_{1}(G)$. 
        Therefore, $G^{\ast} \in \Phi(r,s,\ell,k,m)$. 
        
        Let $p$ and $q$ be nonnegative integers such that $\hat{m} = p \hat{r} + q$ and $q \in [1, \hat{r}]$. 
        Since $\hat{m} > \hat{r} \hat{k}$, we have $p \ge \hat{k}$. 
        We now consider the two cases separately: $q \ge p+1$ and $q \le p$.         

        Suppose that $q \ge p+1$. 
        Then let $\hat{G}$ be the graph obtained from $G^{\ast}$ by removing the edge set $\{u_i w_{p+2} \colon i \in [q-p, q]\}$ and then adding the set $\{u_{\hat{r}+1} w_i \colon i \in [p+1]\}$. 
        It is clear that $|\hat{G}| = |G^{\ast}| = |G|$, and 
        \begin{align*}
            Z_{1}(\hat{G}) - Z_{1}(G^{\ast})
            & = \sum_{i \in [p+2]} \big( d_{\hat{G}}^{2}(w_i) - d_{G}^{2}(w_i) \big) \\
            & \qquad + \sum_{j \in [q-p, q]} \big( d_{\hat{G}}^{2}(u_j) - d_{G}^{2}(u_j) \big) + d_{\hat{G}}^{2}(u_{\hat{r}+1}) - d_{G}^{2}(u_{\hat{r}+1}) \\
            & = (p+1)\big( (\hat{r}+1)^2 - \hat{r}^2 \big) \\
            & \qquad + (q-p-1)^2 - q^2 + (p+1) \big( (p+1)^2 - (p+2)^2 \big) + (p+1)^2 \\
            & = 2(p+1)(\hat{r} - q)
            \ge 0. 
        \end{align*}
        For the set $I = \{u_1, \ldots,u_{\ell}\}$, we have $\delta_{\hat{G}}(I) \ge p+1 \ge \hat{k}+1 =k$. 
        Hence, $\hat{G} \in \mathcal{B}(r,s,\ell,k,m)$. 
        Moreover, it follows from $Z_{1}(\hat{G}) \ge Z_{1}(G^{\ast})$ that $\hat{G} \in \Phi(r,s,\ell,k,m)$. 
        Note that $\hat{G}$ also satisfies $\Delta_{\hat{G}, \mathrm{right}} \ge \Delta_{\hat{G}, \mathrm{left}}$. 
        However, $\Delta_{\hat{G},\mathrm{right}} = d_{\hat{G}}(w_1) = d_{G}(w_1) + 1 > \Delta_{G,\mathrm{right}}$, contradicting the maximality of $\Delta_{G,\mathrm{right}}$. 

        Now suppose that $q \le p$. 
        Similar to the argument above, let $\hat{G}$ be the graph obtained from $G^{\ast}$ by removing the edge set $\{u_i w_{p+2} \colon i \in [q]\}$ and then adding the set $\{u_{\hat{r}+1} w_j \colon j \in [q]\}$. 
        It is clear that $|\hat{G}| = |G^{\ast}| = |G|$, and 
        \begin{align*}
            Z_{1}(\hat{G}) - Z_{1}(G^{\ast})
            & = \sum_{i \in [q]} \big( d_{\hat{G}}^{2}(w_i) - d_{G}^{2}(w_i) \big) + d_{\hat{G}}^{2}(w_{p+2}) - d_{G}^{2}(w_{p+2}) \\
            & \qquad + \sum_{j \in [q-p, q]} \big( d_{\hat{G}}^{2}(u_j) - d_{G}^{2}(u_j) \big) + d_{\hat{G}}^{2}(u_{\hat{r}+1}) - d_{G}^{2}(u_{\hat{r}+1}) \\
            & = q \big( (\hat{r}+1)^2 - \hat{r}^2 \big) - q^2 + q \big( (p+1)^2 - (p+2)^2 \big) + q^2 \\
            & = 2q(\hat{r} - p-1) 
            \ge 0, 
        \end{align*}
        where the last inequality follows from Claim~\ref{CLAIM:dw1>=du1} that $p+1 \le d_{G}(u_1) \le d_{G}(w_1) = \hat{r}$. 
        For the set $I = \{u_1, \ldots,u_{\ell}\}$, we have $\delta_{\hat{G}}(I) \ge p+1 \ge \hat{k}+1 =k$. 
        Hence, $\hat{G} \in \mathcal{B}(r,s,\ell,k,m)$. 
        Moreover, it follows from $Z_{1}(\hat{G}) \ge Z_{1}(G^{\ast})$ that $\hat{G} \in \Phi(r,s,\ell,k,m)$. 
        Notice that $\hat{G}$ also satisfies $\Delta_{\hat{G}, \mathrm{right}} \ge \Delta_{\hat{G}, \mathrm{left}}$. 
        However, since $q \ge 1$, we have $\Delta_{\hat{G},\mathrm{right}} = \hat{r} + 1 > \hat{r} = \Delta_{G,\mathrm{right}}$, contradicting the maximality of $\Delta_{G,\mathrm{right}}$.  
    \end{proof}%CLAIM

\end{document}
